\section{Related work}
\label{sec:related}

We describe each work only as far as its abstract does; the reading of how it
bears on the conjecture is ours.

\paragraph{Length of random network curves (claim P).}
\citet{raghu2017expressive} propose trajectory length, the length of the image
of a one-dimensional input path, as a measure of expressivity, and report that
it grows exponentially with depth. \citet{hanin2022preserve} prove that for ReLU
networks with standard random initialisation the expected length distortion does
not grow with depth. \citet{hanin2019few} bound the average number of activation
patterns of a ReLU network at initialisation in terms of the number of neurons
and the input dimension. By \cref{prop:length}, a curve that does not get longer
cannot get much closer in the nearest-point sense.

\paragraph{Squiggliness after training.}
\citet{dherin2021geometric} argue that stochastic gradient descent implicitly
favours functions of low geometric complexity, which for one-dimensional
regression is the arc length. \citet{bubeck2021universal} show that, for broad
classes of models and data, fitting the data smoothly requires $d$ times more
parameters than merely interpolating it, where $d$ is the input dimension; in
that account parameters buy smoothness, the opposite of (P).
\citet{humayun2024grok} report that the linear regions of a network migrate
away from the training samples late in training, and
\citet{hosseini2023straighten} that trained language models straighten the
trajectory of a sentence across layers while untrained ones do not.

\paragraph{Size helping before or without training.}
\citet{frankle2019lottery} find sparse subnetworks that train in isolation to
the accuracy of the full network. \citet{ramanujan2020hidden} show that a random
network with fixed weights contains subnetworks that perform well untrained,
and that these improve as the network gets wider and deeper.
\citet{larsen2022degrees} find that training within a random subspace succeeds
once its dimension passes a threshold that depends on the target loss and the
initial loss. \citet{chiang2023loss} report that sampling random parameter
vectors until one fits the training set reaches test accuracy comparable to
gradient descent in low-sample regimes. These are statements about subnetworks
or about weight space, not about the nearest-point distance of the full random
function, and \cref{cor:he} is compatible with them.
\citet{martinelli2026escape} argue that the lottery-ticket analogy does not
explain overparameterization and propose that width adds dimensions along which
optimisation escapes bad minima.

\paragraph{Wide networks.} In the lazy regime a model behaves like its
linearisation around the initial weights \citep{chizat2019lazy}, with the neural
tangent kernel governing training \citep{jacot2018ntk}, and wide random networks
converge to Gaussian processes \citep{matthews2018gaussian}. This is the setting
in which ``starts closer'' is true, and \cref{sec:param} is its simplest case:
a model that is exactly linear in its parameters, where the kernel does not
move during training.
